% Part: lambda-calculus
% Chapter: lambda-definability
% Section: lists

\documentclass[../../../include/open-logic-section]{subfiles}

\begin{document}

\olfileid{lam}{ldf}{lst}
\olsection{القوائم}

تُحد القائمة~$[\OLId{latin-looped}{M}_\OLMathNumeral{1}, \OLId{latin-looped}{M}_\OLMathNumeral{2}, \ldots, \OLId{latin-looped}{M}_\OLId{latin-ordinary}{n}]$ على الوجه الآتي:
\[
  [\OLId{latin-looped}{M}_\OLMathNumeral{1}, \OLId{latin-looped}{M}_\OLMathNumeral{2}, \ldots, \OLId{latin-looped}{M}_\OLId{latin-ordinary}{n}] \ident \lambd[sz][\OLId{latin-ordinary}{s} \OLId{latin-looped}{M}_\OLMathNumeral{1} (\OLId{latin-ordinary}{s} \OLId{latin-looped}{M}_\OLMathNumeral{2} (\ldots (\OLId{latin-ordinary}{s} \OLId{latin-looped}{M}_\OLId{latin-ordinary}{n} \OLId{latin-ordinary}{z})))]
\]
وبعبارة غير صورية، الحد الممثل للقائمة هو دالة التراكم فيها: تقبل حدين؛ أحدهما
دالة تأخذ القيمة المتراكمة وعنصراً جديداً فترد قيمة متراكمة جديدة، والآخر
القيمة التي يبدأ بها التراكم. ثم تضم العناصر واحداً بعد واحد، وترد الحاصل
في آخر العمل.

\begin{digress}
  ومن عرف البرمجة الدالية رأى الحد شبيهاً بـ\emph{الطي من اليمين}؛ إذ
  جُعلت القائمة دالة تطوي عناصرها من جهة اليمين.
\end{digress}

ومن هذا الترميز يسهل أن نحد طائفة من الدوال النافعة:
\begin{align*}
  \fn{Sum} & \ident
  \lambd[\OLId{latin-ordinary}{l}][\OLId{latin-ordinary}{l} \, (\lambd[xa][\fn{Add} \, \OLId{latin-ordinary}{x} \, \OLId{latin-ordinary}{a}])\,  \num{\OLMathNumeral{0}}]\\
  \fn{Len} & \ident
  \lambd[\OLId{latin-ordinary}{l}][\OLId{latin-ordinary}{l} \, (\lambd[xa][\fn{Add} \, \OLId{latin-ordinary}{a} \, \num{\OLMathNumeral{1}}]) \num{\OLMathNumeral{0}}]
\end{align*}

تحسب~$\fn{Sum}$ مجموع قائمة من أرقام تشرش؛ فتجمع القائمة مبتدئة بالقيمة
~$\num{\OLMathNumeral{0}}$، وتجعل لكل عنصر~$\fn{Add} \, \OLId{latin-ordinary}{x}\, \OLId{latin-ordinary}{a}$ قيمة التراكم الجديدة،
حيث~$\OLId{latin-ordinary}{x}$ هو العنصر الحاضر و$\OLId{latin-ordinary}{a}$ القيمة المتراكمة. والحاصل مجموع العناصر
كلها.

\begin{prob}
  ما عمل~$\fn{Len}$؟ بيّنه.
\end{prob}

\begin{prob}
  حدّ دالة تقلب ترتيب قائمة.
\end{prob}
\end{document}
